
% Idioma y codificación
\usepackage[spanish,es-tabla]{babel}
\usepackage[utf8]{inputenc}
\usepackage[T1]{fontenc}
\usepackage{mathptmx}
\usepackage{amsmath}
\usepackage{amssymb}
% Tablas bonitas
\usepackage{booktabs}
\usepackage{makecell}
\usepackage{tabularx}
\usepackage[table]{xcolor}
\usepackage{longtable}
\usepackage{array}
\newcolumntype{L}[1]{>{\raggedright\arraybackslash}p{#1}}
\usepackage{amsmath}
\definecolor{headerblue}{RGB}{26,97,133}
\definecolor{lineblue}{RGB}{82,172,220}
\definecolor{rowblue}{RGB}{173,211,230}

\arrayrulecolor{lineblue}
\setlength{\arrayrulewidth}{0.5pt}

\usepackage[ruled,vlined]{algorithm2e}
\newcommand{\headercell}[1]{\textcolor{white}{\textbf{#1}}}
% Control de figuras/tablas y referencias
\usepackage{caption}
\usepackage{subcaption} % subfiguras
\usepackage{tikz}
\usetikzlibrary{arrows.meta, positioning}
% Para que [H] funcione y no se muevan tanto
\usepackage{float}
\usepackage{tikz}
\usepackage{xcolor}
% Links (índice clickeable, referencias clickeables)
\usepackage[colorlinks=true,linkcolor=black,citecolor=black,urlcolor=blue]{hyperref}

% Listas mejoradas
\usepackage{enumitem}
\definecolor{azulESCOM}{RGB}{82,172,220}
% Código (recomendado)
\usepackage{listings}
\usepackage{xcolor}
% Bibliografía IEEE

\usepackage[backend=biber,style=ieee]{biblatex}

% Evita los parentesis vacios "()" en referencias web sin fecha de publicacion
\usepackage{xpatch}
\xpatchbibdriver{online}
  {\printtext[parens]{\usebibmacro{date}}}
  {\iffieldundef{year}{}{\printtext[parens]{\usebibmacro{date}}}}
  {}{}
\addbibresource{references.bib}
\setlength{\parindent}{0pt}%sangria
\setlength{\parskip}{0.4em}

\usepackage{csquotes}

% Gráficos y márgenes
\usepackage{graphicx}
\usepackage[
letterpaper,
left=3cm,
right=2cm,
top=2cm,
bottom=2cm
]{geometry}

% Numeracion de tablas y figuras: la clase book ya las numera por capitulo
% (Tabla 3.1, Figura 3.2, ...), por eso ya no se usa \numberwithin.

% Las paginas en blanco que se agregan para que un capitulo empiece a la
% derecha salen sin encabezado ni numero.
\usepackage{emptypage}

% Pie de pagina: numero de pagina a la izquierda y capitulo a la derecha
% (igual en paginas pares e impares); sin encabezado.
\usepackage{fancyhdr}
\pagestyle{fancy}
\fancyhf{}
\renewcommand{\headrulewidth}{0pt}
\fancyfoot[L]{\thepage}
\fancyfoot[R]{\small\nouppercase{\leftmark}}
% El pie muestra "Capítulo N. Nombre"; los apartados no cambian el pie
\renewcommand{\chaptermark}[1]{\markboth{\chaptername~\thechapter.~#1}{}}
\renewcommand{\sectionmark}[1]{}
% Las paginas donde inicia un capitulo usan el mismo pie
\fancypagestyle{plain}{%
	\fancyhf{}%
	\renewcommand{\headrulewidth}{0pt}%
	\fancyfoot[L]{\thepage}%
	\fancyfoot[R]{\small\nouppercase{\leftmark}}%
}

% Margen de holgura para lineas dificiles de partir (URLs largas en la bibliografia)
\setlength{\emergencystretch}{3em}
